% Analisis de la sensibilidad a los parametros (se incluye desde main.tex).

Los resultados muestran que los parámetros determinan si el ecosistema es estable:
\begin{itemize}[nosep]
  \item \textbf{\texttt{GEN\_FOOD\_FOXES} bajo (6).} Los zorros mueren de hambre antes de encontrar presas
    suficientes. Su población cae de forma sostenida hasta casi extinguirse (43 individuos en la generación
    1500), y los conejos, con cada vez menos depredadores, crecen hasta superar los 10\,000.
  \item \textbf{\texttt{GEN\_FOOD\_FOXES} alto (16).} Los zorros resisten largos periodos sin comer. El
    ciclo depredador--presa es el más visible: en las primeras 300 generaciones se observan varias
    oscilaciones amortiguadas, con los máximos de zorros desfasados respecto de los de conejos. Después,
    ambas especies se estabilizan, con una población de zorros alta (unos 2\,000) y los conejos en un nivel
    más bajo que con los parámetros de referencia.
  \item \textbf{\texttt{GEN\_PROC\_RABBITS} bajo (2).} Los conejos se reproducen muy rápido y sostienen la
    mayor población de zorros de todos los experimentos (unos 3\,000), con fluctuaciones rápidas y amplias
    en ambas especies.
  \item \textbf{\texttt{GEN\_PROC\_RABBITS} alto (9).} Los conejos se reproducen lentamente. Los zorros no
    encuentran alimento suficiente y quedan al borde de la extinción (43 individuos en la generación 1500),
    mientras los conejos, liberados de la depredación, aumentan.
\end{itemize}
En resumen, el ecosistema es estable solo cuando la tasa de reproducción de las presas y la resistencia
al hambre de los depredadores están equilibradas. Si no lo están, los zorros se extinguen. Las rocas actúan
como refugios locales: fragmentan el espacio y dificultan que los zorros alcancen a los conejos, lo que
ayuda a que las presas sobrevivan a los periodos de mayor depredación. Estos experimentos pueden
reproducirse de forma interactiva en el visualizador modificando los parámetros durante la ejecución.
